\documentclass[10pt]{article}
\usepackage[margin=20mm]{geometry}
\usepackage{amsmath,amssymb,fontspec,unicode-math,xurl,hyperref}
\IfFontExistsTF{DejaVu Serif}{\setmainfont{DejaVu Serif}}{\setmainfont{Latin Modern Roman}}
\setmathfont{latinmodern-math.otf}
\setlength{\parindent}{0pt}\setlength{\parskip}{6pt}\sloppy
\begin{document}
Abstract: The printed core diffusion inference algorithm fails an exact Gaussian check even with an optimal noise predictor. The probability-distance definition and bootstrap significance procedure also need correction; historical experiment impact remains unresolved. Independent diagnostics and static inspection distinguish these issues from the substantially more complete current sampler. Historical experiment identity, text-conditioning evaluation and human ranking protocols remain evidence requests.\par
Author: AI-assisted technical review\par
Date:\par
Title: Technical review of BRIDGE\par
\# Identity, scope and outcome\par
Hao Li, Yu-Hao Huang, Chang Xu, Viktor Schlegel, Renhe Jiang, Riza Batista-Navarro, Goran Nenadic and Jiang Bian, BRIDGE: Bootstrapping Text to Control Time-Series Generation via Multi-Agent Iterative Optimization and Diffusion Modeling, ICML 2025, PMLR 267:34742–34773, OpenReview . Main pp. 1–9 and the applicable appendices on pp. 14–32 were read; all seventeen figures were inspected as rendered source pages. References were not independently certified. Anchors use physical PDF pages.\par
The independently checked assessment is Major for the invalid published core inference algorithm, with additional Minor findings and unresolved protocol and historical-version evidence. The paper supplies useful positive generation results, but the printed sampler, metric definitions and bootstrap significance procedure require corrections. Inspection of the public repository is substantial counterevidence to treating the sampler’s documentation error as an implemented historical failure. Neither the original experiment outputs nor their exact execution revision was available; no author program was run.\par
\# P1: Algorithm 2 repeatedly assigns an estimate of the clean sample to the next noisy state\par
Equation reference: eq:printed. Algorithm 1 on p. 20 defines the standard cumulative-noise forward process \[x_n=\sqrt{\bar\alpha_n}x_0+\sqrt{1-\bar\alpha_n}\epsilon.\] Algorithm 2, line 6 on that page, assigns \[
 \widehat x_{n-1}=\frac{\widehat x_n-\sqrt{1-\bar\alpha_n}\widetilde\epsilon_n}{\sqrt{\bar\alpha_n}}.\] This is the usual reconstruction of \(x_0\), not a complete transition to noise level \(n-1\). The loop then treats it as the latter. This is a substantive pseudocode error, rather than an extraction artifact: the formula was checked in the rendered PDF.\par
A small analytic diagnostic takes \(x_0\sim N(0,1)\), a fixed uninformative condition, and the exact square-loss noise predictor \(\epsilon^*(x_n,n)=\sqrt{1-\bar\alpha_n}x_n\). All forward marginals are \(N(0,1)\). Equation [eq:printed] (\#eq:printed) multiplies \(x_n\) by \(\sqrt{\bar\alpha_n}\). Starting from unit Gaussian noise, its final variance is therefore \(\prod_{n=1}^N\bar\alpha_n\). With \(\alpha_n=0.99\) and \(N=100\), it is \(0.99^{5050}=9.07\times10^{-23}\), rather than one. An exact Gaussian reverse transition \(x_{n-1}=\sqrt{\alpha_n}x_n+\sqrt{1-\alpha_n}z\) preserves unit variance. This counterexample concerns the literal printed algorithm under its own noise-prediction interpretation; it is not a simulation of the trained BRIDGE network.\par
The pinned public implementation materially limits the implication. At repository revision dated 2026-10-19, , lines 187–194, computes the clean estimate and then adds the previous-level coefficient, direction term and stochastic term. The DDPM path in , including lines 220–256 and 714–767, also has posterior coefficients and noise. The main generation test helpers (lines 80, 127, 254 and 333) request DDPM sampling; the plot helper at line 200 requests DDIM. Thus the missing transition is present in current code; there is no basis to claim it was absent in the historical runs. Replace Algorithm 2 with the actual sampler, specify the noise schedule and variance choice, and identify the experiment revision. The printed core inference specification is therefore assessed as Major: a necessary generation algorithm fails even with the exact noise predictor. A standard repair and survival of other results do not validate that algorithm. Current-code counterevidence prevents extending this conclusion to the historical implementation or all reported experiments; that separate question remains unresolved.\par
\# P2: The documented MDD and the implemented histogram scores are different quantities\par
Appendix G, pp. 23–24, defines \(\mathrm{MDD}(P,Q)=\sum_x|P(x)-Q(x)|\) for probability distributions. On a discrete probability interpretation, it is bounded by two, whereas Table 1 includes values such as Solar \(375.531\). This does not make the table impossible: the pinned code implements a different, unbounded density-based quantity.\par
In , `HistoLoss` (lines 107–148) forms per-time histograms and averages absolute differences of bin densities, without multiplying by bin width. For two bins with probabilities \((1/2,1/2)\) and \((1,0)\), the probability \(L^1\) distance is one, but the mean density difference for width \(0.001\) is 500. The result depends on units, preprocessing, bins and the aggregation convention. The numerical table can be consistent with this definition. Publish the actual formula and preprocessing, or recompute a probability-distance metric; do not compare the reported magnitude with a theoretical probability-distance bound.\par
There is also a concrete limitation in the current KL estimator. Lines 10–16 construct real-data bin edges, then call `numpy.histogram` with `density=True` for generated samples on those edges and pass the result to `scipy.stats.entropy`. Generated values outside the real range are discarded and the in-range histogram is renormalized. An independent calculation with one percent matching in-range samples and 99 percent samples at 100, while the real data lie in \([0,1]\), gives KL \(8.7\times10^{-17}\). The score can consequently miss arbitrarily large out-of-range mass. This is a property of the inspected estimator, not evidence that the paper’s generated samples actually have those outliers. The MDD implementation need not have the same blind spot because its generated counts use the total sample count.\par
Retain overflow bins or a declared common support, report the excluded fraction, and specify smoothing and normalization. Re-evaluate original outputs to determine whether rankings change. Current-code metric behavior is verified; identity with the historical evaluator and any effect on the paper’s rankings remain unresolved.\par
\# P3: Bootstrap replicates are not additional independent experimental units\par
Appendix M, p. 27, states that OWA is resampled 1000 times and a \(t\)-test is conducted on the resulting subset. If the resampled OWA values are treated as independent observations in an ordinary \(t\)-test, the procedure divides uncertainty by approximately \(\sqrt{1000}\) a second time. Bootstrap replication estimates a sampling distribution; increasing its Monte Carlo resolution does not increase the original number of test series.\par
An independent diagnostic uses 20 paired differences, ten each of \(-0.9\) and \(1.1\). The original paired \(t\)-test has \(p=0.668\). With seed 359, 1000 bootstrap means have percentile interval \([-0.3,0.5]\), while a one-sample \(t\)-test treating those means as the sample gives \(p=7.70\times10^{-38}\). This is a deliberately constructed example, not a reconstruction of Table 14 or a claim that its true conclusions reverse. It shows why the stated procedure cannot by itself validate those significance marks.\par
For OWA, resample the appropriate original series units jointly across methods, recompute the full statistic and its normalizing reference within the declared protocol, and use a valid interval or null-calibrated test. Account for dependence where series are not independent and state the multiple-comparison policy. If the actual test instead used a valid bootstrap tail probability or another correction, provide the implementation and amend the description. The direct SMAPE/MASE tests mentioned on the same page are a separate procedure and are not invalidated by this bootstrap argument. Original paired errors and the test script are required to verify the significance claims.\par
\# Current text-conditioning code and unresolved evaluation details\par
The available exposes a `description` argument, but both `\_predict\_llmtime` (lines 102–140) and `\_predict\_promptcast` (143–182) never read it. Static AST inspection independently confirms zero loads in either function. Thus this wrapper, with all other inputs fixed, cannot measure a prediction change caused by changing that text argument. This is a verified local implementation limitation; it does not establish that this wrapper produced Table 5. The inspected entry point also permits a different predictor arrangement. Supply the original text-conditioning call path before attributing the reported refinement improvements to this public wrapper; wire the argument into the model request if that wrapper is intended to reproduce them.\par
The paper explicitly distinguishes optimizing description templates on AirPassengers/Sunspots from filling the frozen templates with per-instance statistics in other domains. Conditioning a time-series generator on statistics of the desired target is part of the stated task. It is not, by itself, forecasting leakage. However, the separately described zero-shot forecasting evaluation and its future-statistic examples on pp. 15–16 require an information-cutoff account: which target-window statistics were available to the prediction prompt, and which windows were reserved for evaluation after template refinement. The current confirms per-window statistic filling, but cannot establish the historical split or the causal source of Table 5 gains. Conditional generation, held-out-domain transfer and forecasting must each state their information sets.\par
Human ranking also needs its exact candidate pool. Table 3 gives HE-Rank 4.2 for the no-text variant on Traffic and Taxi. If the pool is the three displayed variants plus ground truth, a mean ordinal rank cannot exceed four. Appendix H says “different model settings” and does not enumerate the pool, so additional unlisted candidates or a score-versus-rank convention could explain 4.2. This is an evidence request, not proof of invalid annotations. Provide candidate counts, rubric, number of annotators and items, raw rankings and aggregation; distinguish HE-Rank from HE-Mixed, which adds three distractors. The prose claim of an improvement over three points for Traffic and Pedestrian also does not follow from the displayed HE-Rank differences of 2.8 and 1.1.\par
\# Positive checks, coverage and resolution\par
The full model improves JFTSD over the no-text variant in all twelve displayed domains, and beats both ablations in eleven; Rain has a slightly better no-prototype value. Against external baselines, Table 1 gives the best MDD in ten domains and best KL in eleven, so a count of eleven must identify the metric. Table 5’s four initial-to-refined relative improvements are \(17.06\%\), \(18.06\%\), \(18.77\%\) and \(14.77\%\); these support the general improvement direction, with the last slightly below the prose’s \(15\%\). These arithmetic checks neither substitute for uncertainty estimates nor imply that domain averages are independent replicates.\par
All source figures, including the prototype panels on pp. 28–32, were viewed. They illustrate conditional sample behavior but cannot establish distributional calibration from selected examples. Main text, Algorithms 1–2, description construction, model and dataset specifications, metrics, human evaluation, ablations and downstream evaluation were inspected. The selected public files cover text construction and refinement, evaluation metrics, the prototype module and the relevant diffusion sampling/loss paths. Eight additional files were retrieved through the authorized repository connector at the same pinned revision, with exact Git-blob and SHA-256 checks recorded. Large diffusion and utility files were inspected at the relevant functions, not certified in full. No original model training, author code execution, external literature replication, complete dataset audit, annotation audit or historical checkpoint verification occurred. Separately hosted OpenReview attachments were unavailable during source preparation; embedded appendices were available.\par
Correct the printed algorithm and metric specification, replace or clarify the OWA significance procedure, and connect the reproduction package to the historical run. The latter evidence requests are unresolved. Author adoption is unknown, and repairs beyond the demonstrated algebra and diagnostic examples have not been independently verified.\par
2 H. Li et al. BRIDGE: Bootstrapping Text to Control Time-Series Generation via Multi-Agent Iterative Optimization and Diffusion Modeling. PMLR 267 (2025), 34742–34773. https://proceedings.mlr.press/v267/li25ah.html. Microsoft, TimeCraft public repository, revision . https://github.com/microsoft/TimeCraft. Current public version; historical experiment identity unestablished.\par
\# Reproducibility and roles\par
The matching directory contains the independent check script and its actual output, source and code hashes, coverage, findings, compiler logs and page-level visual records. Author programs were not executed. The source PDF SHA-256 is . Separate assigned GPT-6 Astra / Max units performed the initial review and the subsequent independent source-based check recorded in the verification files. Scientific adjudication remains pending; human reviews are empty.\par
\end{document}
